\chapter{Why Mobile Application Development Frameworks are needed}\label{chap:why_needed}

Native development serves as the baseline for mobile applications, utilizing first-party languages and frameworks like Swift for iOS and Kotlin for Android. While native development provides peak performance, fluid user experiences, and immediate access to new operating system features, it introduces a significant challenge: the necessity to maintain two entirely separate codebases. This dual maintenance requires hiring specialized engineering talent for each platform, which significantly increases development time and overall expenses.

Mobile application development frameworks emerged to solve these exact inefficiencies. They offer alternative architectural paradigms, such as WebView-driven designs, bridge-based Java-Script runtimes, or compiled custom rendering engines, to allow developers to share code across multiple platforms. By utilizing these frameworks, development teams can carefully balance crucial metrics like performance, security, and developer experience (DX) while reducing the friction of context switching between different native languages. 

Ultimately, modern frameworks aim to bridge the gap between high-performance native execution and the efficiency of a shared codebase. This enables rapid development and cost-effectiveness without forcing users to sacrifice the familiar look and feel of a purely native application.